
% Glossar

\chapter*{Glossar}

\section{Java}

\begin{description}
	\item[Klasse] Sammlung an Methoden und Variablen bzw. Definition einer Datenstruktur.
	\item[Konstruktor] Spezielle Methode einer Klasse, die beim instantiieren dieser aufgerufen wird.
	\item[Kindklasse, erweitern] Eine neue Klasse, die Kindklasse, übernimmt alle (mindestens
		als \texttt{protected} deklarierte) Methoden und Variablen der erweiterten Klasse,
		der Elternklasse.
	\item[\texttt{abstract}] Java-Schlüsselwort. Bei Klassen: Klasse kann nicht direkt
		instantiiert werden. Bei Methoden: Kindklassen \textbf{müssen} diese überschreiben.
	\item[Interface] Spezielle Art von Klasse, welche ausschließlich abstrakte Methoden besitzt.
	\item[Enumeration, Enum]: Spezielle Art von Klasse, welche eine Feste Anzahl an benannten
		Instanzen definiert, also nicht beliebig instantiiert werden kann.
	\item[\texttt{public}] Java-Schlüsselwort. Klasse, Methode oder Variable kann von
		überall erreicht werden.
	\item[\texttt{protected}] Java-Schlüsselwort. Klasse, Methode oder Variable kann
		nur innerhalb des gleichen Paketes oder von Kindklassen erreicht werden.
	\item[\texttt{private}] Java-Schlüsselwort. Klasse, Methode oder Variable kann
		nur innerhalb ihrer Klasse erreicht werden.
	\item[\texttt{static}] Java-Schlüsselwort. Klasse, Methode oder Variable kann
		unabhängig von einer Instanz verwendet werden, gehört also direkt zu einer
		Klasse.
	\item[Map, HashMap, LinkedHashMap] Eindeutige Zuordnung von Wertepaaren.
	\item[Schlüsselmenge] Menge der Werte, denen in einer \texttt{Map} ein Wert zugeordnet ist.
\end{description}


\section{Spiel}

\begin{description}
	\item[Spiel-Engine, Engine] Code-Grundgerüst bzw. Bibliothek, welche maßgebliche
		Unterstützung beim Erstellen eines Spiels bietet.
	\item[Spiel-Objekt, Objekt] Akteur innerhalb des Spiels. Hat eine Position, Größe und optional ein Bild.
	\item[Szene] Eine Gruppe von Objekten.
\end{description}